% SPDX-License-Identifier: CC-BY-NC-ND-4.0
% Copyright 2026 Ingolf Lohmann.
\documentclass[11pt,a4paper]{article}

\usepackage{fontspec}
\setmainfont{Latin Modern Roman}
\setsansfont{Latin Modern Sans}
\setmonofont{Latin Modern Mono}[Scale=0.86]
\usepackage[provide=*,main=ngerman,english]{babel}
\usepackage{microtype}
\usepackage{geometry}
\usepackage{amsmath,amssymb,mathtools,amsthm}
\usepackage{booktabs,longtable,tabularx,array}
\usepackage{enumitem}
\usepackage{xcolor}
\usepackage{fancyhdr}
\usepackage{seqsplit}
\usepackage{xurl}
\usepackage{hyperref}
\usepackage{bookmark}
\usepackage[nameinlink,noabbrev]{cleveref}
\newcolumntype{L}[1]{>{\raggedright\arraybackslash}p{#1}}

\geometry{left=24mm,right=24mm,top=24mm,bottom=27mm,headheight=14pt}
\setlength{\parindent}{0pt}
\setlength{\parskip}{0.55em}
\setlength{\emergencystretch}{3em}
\setlist[itemize]{leftmargin=1.45em,itemsep=0.18em,topsep=0.3em}
\setlist[enumerate]{leftmargin=1.7em,itemsep=0.22em,topsep=0.3em}

\definecolor{QBlue}{HTML}{153E69}
\definecolor{QLight}{HTML}{ECF3F8}
\definecolor{QGreen}{HTML}{E8F3EA}
\definecolor{QAmber}{HTML}{FFF2CC}
\definecolor{QRed}{HTML}{FBE9E7}
\definecolor{QGray}{HTML}{4D5660}

\hypersetup{
  colorlinks=true,
  linkcolor=QBlue,
  citecolor=QBlue,
  urlcolor=QBlue,
  pdftitle={QIK-VRT: Epistemische Fairness, Vorstellungskraft und der Weg vom virtuellen Zeugen zum Beobachtungssystem},
  pdfauthor={Ingolf Lohmann},
  pdfsubject={Wissenschaftlicher Nachtrag zum bidirektionalen virtuellen Zeitkanal},
  pdfkeywords={QIK-VRT, virtuelle Zeit, Informationsübertragung, Beobachtungssystem, epistemische Fairness, Vorstellungskraft, Retrokausalität}
}

\pagestyle{fancy}
\fancyhf{}
\fancyhead[L]{\small\textcolor{QGray}{QIK--VRT: wissenschaftlicher Nachtrag}}
\fancyhead[R]{\small\textcolor{QGray}{Ingolf Lohmann}}
\fancyfoot[C]{\small\thepage}
\renewcommand{\headrulewidth}{0.3pt}

\newtheorem{definition}{Definition}[section]
\newtheorem{satz}{Satz}[section]
\newtheorem{korollar}{Korollar}[section]
\newtheorem{lemma}{Lemma}[section]
\theoremstyle{remark}
\newtheorem{bemerkung}{Bemerkung}[section]
\crefname{definition}{Definition}{Definitionen}
\crefname{satz}{Satz}{Sätze}
\crefname{korollar}{Korollar}{Korollare}
\crefname{lemma}{Lemma}{Lemmata}

\newcommand{\qikvrt}{QIK--VRT}
\newcommand{\DONE}{\texttt{EFFECT\_ACK\_DONE}}
\newcommand{\hash}[1]{{\ttfamily\small\seqsplit{#1}}}
\newcommand{\statusbox}[2]{%
  \par\medskip\noindent
  \colorbox{#1}{\parbox{\dimexpr\linewidth-2\fboxsep\relax}{#2}}%
  \par\medskip
}

\begin{document}
\special{pdf:trailerid [<059d4c78efd923194210801077d7eb45><059d4c78efd923194210801077d7eb45>]}

\hypersetup{pageanchor=false}
\begin{titlepage}
  \thispagestyle{empty}
  \vspace*{10mm}
  {\color{QBlue}\rule{\textwidth}{1.4pt}}\\[1.0em]
  {\Huge\bfseries QIK--VRT\par}
  \vspace{0.28em}
  {\LARGE\bfseries Epistemische Fairness und Vorstellungskraft\\[0.12em]
  Vom virtuellen Zeugen\\[0.12em]
  zum Beobachtungssystem\par}
  \vspace{0.9em}
  {\LARGE Wissenschaftlicher Nachtrag zum bidirektionalen virtuellen Zeitkanal\par}
  \vspace{1.0em}
  {\color{QBlue}\rule{\textwidth}{0.6pt}}

  \vfill
  {\large\textbf{Ingolf Lohmann}\par}
  \vspace{0.35em}
  {\normalsize Independent Researcher; QIK--VRT\par}
  \vspace{0.35em}
  {\normalsize Addendum 1.0.0 zum Kandidaten 1.0.0 -- 1. August 2026\par}

  \vfill
  \statusbox{QLight}{
    \textbf{Additiver Befund.}
    Der ausgeführte endliche QIK--VRT-Zeuge bleibt unverändert: Er überträgt
    vollständige endliche Nachrichten in beiden Richtungen der virtuellen
    Zeitadressierung, während die Hostzeit strikt vorwärts läuft. Dieser
    Nachtrag zeigt, wie daraus überprüfbare Geräte- und
    Beobachtungshypothesen gebildet werden können. Er definiert ein
    Beobachtungsprädikat, eine Realisierungsleiter und Schutzregeln gegen
    epistemische Übervorteilung. Eine reale oder physikalische Brücke wird
    dadurch untersuchbar, aber nicht ohne Messung als vorhanden ausgegeben.
  }

  \vfill
  {\small
    Dokumentationslizenz: CC BY-NC-ND 4.0.\\
    Vorveröffentlichungskandidat; nicht peer-reviewed; noch kein neuer DOI.\\
    Der begleitende Internet-Draft-Entwurf ist nicht beim IETF eingereicht
    und bezeichnet keinen IETF-Konsens.
  }
\end{titlepage}
\hypersetup{pageanchor=true}
\setcounter{page}{1}

\selectlanguage{ngerman}

\begin{abstract}
Dieser Nachtrag setzt den bytegebundenen QIK--VRT-Kandidaten zum
bidirektionalen virtuellen Zeitkanal voraus und ändert keinen seiner dreizehn
Claims. Der vorhandene ISO-C90-Zeuge realisiert einen endlichen Dialog: Eine
257-Byte-Anfrage wird von der virtuellen Adresse 30 an die frühere Adresse 15
übertragen; dort wird sie deterministisch verarbeitet; eine 258-Byte-Antwort
kehrt zur Adresse 30 zurück. Sämtliche erzeugenden Hostereignisse bleiben
streng vorwärtsgeordnet. Unter expliziten Endlichkeits-, Integritäts- und
Zustellannahmen ist die Konstruktion auf beliebige einzelne endliche
Nachrichten verallgemeinert. Dieser Informatikbefund wird hier nicht neu
bewiesen, sondern exakt als Stamm des erweiterten Erkenntnisbaums gebunden.

Zwei neue Audioquellen motivieren die Anschlussfragen, welche Rolle
Vorstellungskraft bei der Bildung technischer Hypothesen spielt, wann ein
Computerspiel oder anderes Softwaresystem zu einem Beobachtungssystem wird
und wie verdeckte Wissens- oder Zugriffsvorsprünge begrenzt werden müssen.
Wir führen eine achtstufige Realisierungsleiter von der Idee bis zur
autorisierten Außenwirkung ein und definieren ein Beobachtungssystem durch
Datenerfassung, Kalibrierung, Zeitbindung, Provenienz, Integrität,
Falsifizierbarkeit und Governance. Ein einfacher Gegenmodellsatz zeigt, warum
Softwaremöglichkeit allein keinen entsprechenden Naturvorgang impliziert.
Gleichzeitig wird konstruktiv präzisiert, unter welchen Bedingungen
instrumentierte Software tatsächlich als Beobachtungssystem fungieren kann.

Epistemische Fairness wird als normative Schutzschicht operationalisiert:
Simulation, Antizipation, Beobachtung, Verifikation und autorisierte Wirkung
müssen unterscheidbar bleiben; privilegierter Zugang darf Provenienz,
Anfechtbarkeit oder Einwilligung nicht ersetzen. Behauptungen über bereits
verdeckte reale Systeme bleiben offen und testbedürftig. Ebenso werden weder
Superdeterminismus noch eine physikalische Zukunft-zu-Vergangenheit-
Übertragung behauptet. Ein begleitendes, individuelles Experimental-
Anwendungsprofil übersetzt die überprüfbaren Teile in Anforderungen über dem
unveränderten EFFECT\_ACK-Wire-v1-Datensatz.
\end{abstract}

\begin{otherlanguage}{english}
\begin{abstract}
This addendum is bound to the byte-identical QIK--VRT prepublication
candidate on the bidirectional virtual-time channel and does not rewrite any
of its thirteen claims. The existing ISO C90 witness executes a finite
dialogue: a 257-byte request travels from virtual address 30 to earlier
address 15, is processed deterministically, and returns a 258-byte response
to address 30, while all generating host events remain strictly forward
ordered. Under explicit finiteness, integrity, and delivery assumptions, the
construction applies to every individual finite message. That computer
science result is inherited here as the fixed trunk of an extended evidence
tree.

Two new audio sources motivate the next questions: how imagination can form
technical hypotheses, when a game or other software becomes an observation
system, and how hidden informational or access advantages should be
constrained. We introduce an eight-level realization ladder from articulated
idea to authorized external effect. We define an observation system through
acquisition, calibration, temporal binding, provenance, integrity,
falsifiability, and governance. A countermodel establishes that software
possibility alone does not entail a corresponding physical phenomenon;
conversely, explicit instrumentation conditions make the transition from
software to a qualified observation system testable.

Epistemic fairness becomes a normative guard layer: simulation,
anticipation, observation, verification, and authorized effect must remain
distinguishable, and privileged access must not replace provenance,
contestability, or consent. Claims of already hidden real-world deployments
remain empirically open. The addendum proves neither superdeterminism nor
physical future-to-past communication. A companion individual Experimental
application-profile draft translates the mechanically checkable requirements
onto the unchanged EFFECT\_ACK version-1 wire record.
\end{abstract}
\end{otherlanguage}

\setcounter{tocdepth}{2}
\tableofcontents
\newpage

\section{Exakter Ausgangspunkt: Der Erkenntnisbaum wird erweitert}

Der Nachtrag beginnt bewusst nicht wieder bei null. Sein unmittelbarer
Parent ist der Vorveröffentlichungskandidat im Pull Request 293 des
Authority-Repositories \texttt{Goldkelch/qik-vrt}:

\begin{longtable}{@{}L{0.26\linewidth}L{0.66\linewidth}@{}}
\toprule
\textbf{Bindung} & \textbf{Exakter Wert} \\
\midrule
\endhead
Parent-Commit &
\hash{5df3e24496afbeac60dfc78ffb12d673f163ee04} \\
Parent-Tree &
\hash{c147d82b61efc989f0cc0aa698e16bf71c6ec9da} \\
Basis-Publikations-ID &
\texttt{qikvrt-bidirectional-virtual-time-channel-v1} \\
Basis-Machine-Proof &
\hash{ad30282f9ecc30414c5dd7eef0460a9766f03c8184c01166f15dc9cb567aa72a} \\
Delta-Semantik &
\texttt{APPEND\_ONLY\_NO\_BASE\_CLAIM\_REWRITE} \\
\bottomrule
\end{longtable}

Die maschinenlesbare Datei \texttt{INHERITED\_PROOF\_BINDING.json} bindet
alle dreizehn geerbten Claims an ihre unveränderten Klassifikationen und
Statuswerte. Die acht neuen Claims tragen eigene Kennungen
\texttt{VTI-ADD-001} bis \texttt{VTI-ADD-008}. So bleibt nachvollziehbar,
welcher Satz bereits zum Basiswerk gehörte und welcher erst hier hinzukommt.

\statusbox{QGreen}{
  \textbf{Geerbter tragender Informatikbefund.}
  Vollständige einzelne endliche Nachrichten können unter den im Basiswerk
  genannten Bedingungen bytegenau in beide Richtungen einer virtuellen
  Adressordnung transportiert werden. Der konkrete C90-Lauf führt Anfrage
  und Antwort aus; alle tatsächlichen Hostereignisse bleiben strikt
  vorwärtsgerichtet; die gebundene Quellhistorie wird nicht überschrieben.
}

Dieser Befund ist stark, aber genau abgegrenzt. \emph{Bidirektional} bezieht
sich hier auf zwei Transportrichtungen zwischen virtuellen Adressen.
\emph{Vollständig} bedeutet Bytevollständigkeit einer jeweils endlichen
Nachricht innerhalb der spezifizierten Ressourcen und Annahmen. Keines der
beiden Wörter behauptet automatisch einen physikalischen Kanal aus einer
späteren in eine frühere Raumzeitregion.

\subsection{Was unverändert gilt}

\begin{itemize}
  \item \texttt{VTI-001} und \texttt{VTI-002} bleiben ausgeführte Evidenz
    für den endlichen C90-Zeugen und seine benannten Grenzfälle.
  \item \texttt{VTI-007} bindet den vorhandenen Lean-Kernelreceipt nur in
    dessen ausgewiesenem Scope; der Nachtrag etikettiert keinen neuen Satz
    nachträglich als kernelverifiziert.
  \item \texttt{VTI-010} trennt Byteidentität, Bedeutung, Wahrheit und
    autorisierte Wirkung. Diese Trennung wird hier auf Beobachtungsdaten
    ausgedehnt.
  \item Die physikalische Brücke und physikalische Rückwärtssignalisierung
    bleiben offen. Der Status wird weder angehoben noch rhetorisch verwischt.
\end{itemize}

\section{Zwei neue Quellen und ihre wissenschaftliche Rolle}

Die Dateien \emph{Das ist Übervorteilen!} und \emph{Das ist
Vorstellungskraft!} sind neue Primärquellen. Sie wurden lokal und offline
transkribiert. Das Repository enthält weder die Audioinhalte noch eine
behauptete Identitätsauthentifizierung, sondern SHA-256-Medienbindungen,
technische Metadaten, unveränderte ASR-Rohfassungen, konservativ geprüfte
Lesefassungen und ein Unsicherheitsprotokoll.

\begin{longtable}{@{}L{0.24\linewidth}L{0.68\linewidth}@{}}
\toprule
\textbf{Quelle} & \textbf{SHA-256 der kanonischen Audiodatei} \\
\midrule
\endhead
Übervorteilen &
\hash{c7ae25dc1a689211fe9caa60d39cbd2bea3265aab655b4a7fc14daebc1582a05} \\
Vorstellungskraft &
\hash{a4a9d1141c33848b3ee6ef30d030b176b4b888103a09ec14503f65b0b21e19ca} \\
\bottomrule
\end{longtable}

Aus einer Aufnahme folgt zunächst nur, dass eine Aussage in den gebundenen
Bytes enthalten ist. Dateiname, Sprecherzuschreibung, Transkription und
sachliche Wahrheit sind getrennte Ebenen. Der Nachtrag übernimmt deshalb
nicht jede Aussage als Fakt. Er gewinnt aus den Quellen vielmehr zwei
wissenschaftlich bearbeitbare Richtungen:

\begin{enumerate}
  \item \textbf{Vorstellungskraft als Hypothesengenerator:} Welche heute
    verfügbaren Geräte, Datenpfade und Softwareeigenschaften könnten eine
    klar formulierte Folgekonfiguration realisieren?
  \item \textbf{Epistemische Fairness als Wirkungsgrenze:} Welche
    Offenlegungs-, Prüf- und Widerspruchsrechte sind nötig, wenn Wissen oder
    Beobachtungszugänge asymmetrisch verteilt sind?
\end{enumerate}

Die in einer Aufnahme geäußerte These, entsprechende Systeme seien bereits
verdeckt real im Einsatz, ist damit als Autorenaussage quellengebunden. Ohne
benannte Systeme, inspizierbare Artefakte, Messprotokolle und unabhängige
Replikation bleibt ihr Sachgehalt \texttt{OPEN\_EMPIRICAL}.

\section{Die Realisierungsleiter}

\begin{definition}[Realisierungsleiter]
Die Realisierungsleiter ist die geordnete Folge
\[
R_0 \prec R_1 \prec R_2 \prec R_3 \prec R_4 \prec R_5 \prec R_6 \prec R_7,
\]
wobei jede Stufe eine eigene Evidenzpflicht besitzt. Das Erreichen einer
Stufe macht die nächste untersuchbar, ersetzt deren Evidenz aber nicht.
\end{definition}

\begin{longtable}{@{}L{0.08\linewidth}L{0.27\linewidth}L{0.55\linewidth}@{}}
\toprule
\textbf{Stufe} & \textbf{Gegenstand} & \textbf{Benötigter Nachweis} \\
\midrule
\endhead
$R_0$ & artikulierte Idee & verständliche, quellengebundene und abgrenzbare These \\
$R_1$ & konsistentes Softwaremodell & definierte Zustände, Übergänge, Annahmen und Grenzen \\
$R_2$ & ausführbarer virtueller Zeuge & reproduzierbarer Lauf sowie positive und negative Tests \\
$R_3$ & reale Geräteimplementierung & gebundene Hardware, Software, Versionen und Datenpfade \\
$R_4$ & qualifiziertes Beobachtungssystem & vollständiges Beobachtungsprädikat aus \cref{sec:obs} \\
$R_5$ & physikalischer Befund & Kalibrierung, Kontrollen, Messunsicherheit und Replikation \\
$R_6$ & kausal zugerechnete Wirkung & Intervention oder tragfähige Identifikationsstrategie \\
$R_7$ & autorisierte Außenwirkung & Rechte, Zweckbindung, Einwilligung und Wirkungsgate \\
\bottomrule
\end{longtable}

Das Basiswerk erreicht $R_2$ für seinen deklarierten virtuellen Scope. Dass
der Zeuge auf realer Hardware läuft, ist selbstverständlich ein realer
Rechenvorgang. Es belegt jedoch nicht bereits die spezielle physikalische
Eigenschaft, die ein behaupteter Zukunft-zu-Vergangenheit-Kanal benötigen
würde. Dazu wären zusätzliche Übergangsnachweise von $R_3$ bis $R_6$ nötig.

\begin{satz}[Nicht-Implikation von Software und Natur]
Die Existenz oder Ausführung eines Softwaremodells $S$ impliziert für sich
allein nicht die Existenz eines von $S$ repräsentierten physikalischen
Phänomens $P$.
\end{satz}

\begin{proof}
Betrachte zwei mögliche Welten mit derselben Rechenmaschine und demselben
ausgeführten Programm $S$. In Welt $W_1$ existiert zusätzlich der behauptete
physikalische Kanal $P$; in Welt $W_0$ existiert er nicht. Solange $S$ keine
kalibrierte, ausschließende Messbrücke zur Unterscheidung beider Welten
enthält, ist die beobachtete Softwareausführung in $W_0$ und $W_1$ gleich.
Damit ist $S$ mit $\neg P$ vereinbar. Folglich gilt logisch nicht
$S\Rightarrow P$. Für die Ableitung von $P$ ist mindestens eine zusätzliche
Brückenprämisse mit empirischer Evidenz erforderlich.
\end{proof}

Der Satz entwertet das Softwareergebnis nicht. Er zeigt präzise, welche
Arbeit zwischen einem richtigen Informatiksatz und einem neuen Physikbefund
liegt. Genau diese explizite Schnittstelle macht den nächsten Schritt
planbar.

\section{Vom Computerspiel zum Beobachtungssystem}
\label{sec:obs}

Ein Computerspiel ist zunächst ein interaktives Softwaresystem. Es kann
Eingaben erfassen, Zustände simulieren und Telemetrie speichern. Diese
Fähigkeiten allein machen jede Ausgabe noch nicht zu einer belastbaren
Beobachtung über die Außenwelt.

\begin{definition}[Beobachtungsprädikat]
Für einen deklarierten Scope heißt ein System $X$ hier genau dann
\emph{qualifiziertes Beobachtungssystem}, wenn
\[
\operatorname{OBS}(X)=A\land K\land T\land Q\land I\land F\land G
\]
erfüllt ist. Dabei stehen $A$ für spezifizierte Erfassung, $K$ für
Kalibrierung, $T$ für Zeit- und Ereignisbindung, $Q$ für Quellenprovenienz,
$I$ für Daten- und Auswertungsintegrität, $F$ für Falsifizierbarkeit und
$G$ für Governance einschließlich Zweck, Rechten und Auditierbarkeit.
\end{definition}

\begin{satz}[Bedingter Instrumentierungssatz]
Ein Spiel oder anderes interaktives Softwaresystem kann innerhalb eines
genau benannten Scopes als Beobachtungssystem dienen, wenn alle sieben
Konjunkte von $\operatorname{OBS}$ nachweislich erfüllt sind. Aus der bloßen
Systembezeichnung, Interaktion oder Telemetrie folgt diese Eigenschaft nicht.
\end{satz}

\begin{proof}
Die erste Aussage folgt konstruktiv aus der Definition: Man ergänzt einen
spezifizierten Sensor- oder Eingabepfad, bindet dessen Kalibrierung und Zeit,
erhält eine überprüfbare Transformationskette, schützt Rohdaten und Analyse,
präregistriert widerlegbare Ausgaben und setzt Rechte sowie Auditverfahren
durch. Dann sind alle Konjunkte erfüllt. Für die zweite Aussage genügt ein
Gegenbeispiel: Ein Spiel speichert beliebige Klickzahlen, ohne Sensorbezug,
Kalibrierung oder Falsifikationsregel. Telemetrie liegt vor, aber mindestens
$K$, $Q$ und $F$ fehlen; also gilt $\neg\operatorname{OBS}(X)$.
\end{proof}

\subsection{Vier greifbare Beispiele}

\textbf{Beispiel 1: Eingabetelemetrie.}
Ein Spiel zeichnet Tastendrücke auf. Es beobachtet damit zunächst nur
technisch definierte Eingabeereignisse innerhalb seines eigenen Pfads. Eine
Diagnose über Aufmerksamkeit, Absicht oder Gesundheitszustand folgt daraus
nicht ohne validiertes Messmodell, Unsicherheitsangabe und geeignete
Einwilligung.

\textbf{Beispiel 2: Kalibrierter Bewegungssensor.}
Ein Spiel nutzt einen gebundenen Sensor, um einen Rehabilitationsablauf zu
erfassen. Mit dokumentierter Kalibrierung, Zeitbindung, Rohdatenintegrität,
validierter Auswertung und zweckgebundener Einwilligung kann es in diesem
engen Scope ein Beobachtungssystem sein. Es wird dadurch nicht zum
universellen Beobachter.

\textbf{Beispiel 3: Virtueller Replay.}
Ein deterministischer Replay erzeugt für denselben gebundenen Eingabestand
dieselbe Ausgabe. Das ist starke Evidenz für Reproduzierbarkeit innerhalb des
Modells. Es ist keine neue Direktbeobachtung der Außenwelt und kein Nachweis,
dass der simulierte Vorgang physikalisch stattfand.

\textbf{Beispiel 4: Behauptetes verborgenes System.}
Die Aussage, ein unbenanntes Spiel sei heimlich ein umfassendes
Observationssystem, ist eine prüfbare Hypothese. Erforderlich wären mindestens
identifizierbare Binärdateien oder Geräte, Netzwerk- und Sensorpfade,
Erfassungsprotokolle, Testkonten, Kontrollen, rechtlich zulässige Prüfung und
unabhängige Replikation. Ohne diese Belege bleibt die These offen.

\section{Vorstellungskraft als kontrollierter Suchoperator}

Vorstellungskraft ist hier weder Messinstrument noch Wahrheitsoracle. Sie ist
ein produktiver Operator über dem Hypothesenraum. Er kann bekannte Bausteine
neu kombinieren, überraschende Tests vorschlagen und blinde Flecken sichtbar
machen. Damit dieser Vorteil wissenschaftlich nutzbar wird, durchläuft jeder
Kandidat fünf Schritte:

\begin{enumerate}
  \item \textbf{Inventar:} Geräte, Schnittstellen, Protokolle und
    Abhängigkeiten werden benannt.
  \item \textbf{Konstruktion:} Eine neue Konfiguration wird so genau
    beschrieben, dass eine andere Person sie nachbauen könnte.
  \item \textbf{Vorhersage:} Erwartete Beobachtungen und Gegenbeobachtungen
    werden vor Einsicht in neue Messdaten festgelegt.
  \item \textbf{Branch-Trennung:} Simulation, Antizipation und Messarchiv
    bleiben technisch und semantisch unterscheidbar.
  \item \textbf{Statusübergang:} Erst Messung, Provenienzprüfung und
    Autorisierung dürfen den Status erhöhen.
\end{enumerate}

Der QIK--VRT-Kanal ist dafür besonders anschlussfähig: Virtuelle Adresse,
Branch-Replay, Antwortbindung und Wirkungsgate liefern Bausteine, um
Kontrafaktisches zu berechnen, ohne es nachträglich als unveränderte
Vergangenheit auszugeben. Die Vorstellungskraft erweitert die Menge der
prüfbaren Branches; sie überschreibt nicht den Wahrheitsstatus des Archivs.

\section{Epistemische Fairness und Übervorteilung}

\begin{definition}[Epistemische Übervorteilung]
Epistemische Übervorteilung bezeichnet in diesem Nachtrag eine normative
Risikoklasse: Ein Informations-, Beobachtungs- oder Prognosevorsprung wird so
eingesetzt, dass betroffene Parteien Herkunft, Unsicherheit, Vergleichsmaßstab
oder Wirkung nicht angemessen prüfen, anfechten oder autorisieren können.
\end{definition}

Der Begriff beweist nicht, dass in einem konkreten Fall Unrecht geschehen
ist. Dazu wären Kriterien, Belege, Zuständigkeit und ein faires Verfahren
nötig. Er macht jedoch ein technisches Gestaltungsproblem sichtbar. Ein
System kann Bytes korrekt verarbeiten und trotzdem unfair wirken, etwa wenn
es

\begin{itemize}
  \item privilegierte Beobachtungs- oder Prognosezugänge verschweigt,
  \item simulierte, antizipierte und beobachtete Daten vermischt,
  \item Unsicherheit oder Fehlerkosten einseitig verteilt,
  \item Gegenprüfung und Berichtigung verhindert oder
  \item irreversible Wirkungen ohne operationsspezifische Autorisierung
    auslöst.
\end{itemize}

\begin{satz}[Vorsprung ersetzt keine Wirkungserlaubnis]
Innerhalb des hier deklarierten Governance-Modells kann ein Wissens- oder
Zugriffsvorsprung weder den Wahrheitsstatus eines Datums noch die
Autorisierung einer irreversiblen Außenwirkung ersetzen.
\end{satz}

\begin{proof}
Wahrheitsstatus und Wirkungserlaubnis werden als eigene Gate-Prädikate
definiert. Ein Zugriffsvorsprung ist kein Bestandteil ihrer hinreichenden
Bedingungen. Fehlt ein benötigtes Prädikat, bleibt die ordinary release
unabhängig vom Vorsprung falsch. Die Aussage ist eine überprüfbare
Eigenschaft dieser Architektur, kein universelles moralisches Theorem.
\end{proof}

Praktisch benötigt jedes wirkungsrelevante Informationsobjekt mindestens
einen sichtbaren epistemischen Status:

\begin{longtable}{@{}L{0.24\linewidth}L{0.68\linewidth}@{}}
\toprule
\textbf{Status} & \textbf{Bedeutung} \\
\midrule
\endhead
\texttt{CANDIDATE} & formulierter, noch nicht geprüfter Vorschlag \\
\texttt{ANTICIPATED} & vorab erwartetes Ergebnis, keine Beobachtung \\
\texttt{OBSERVED} & erfasster Befund mit deklariertem Mess- und Quellenscope \\
\texttt{VERIFIED} & nach benanntem Verfahren überprüfter Befund \\
\texttt{AUTHORIZED\_EFFECT} & gesondert freigegebene konkrete Wirkung \\
\bottomrule
\end{longtable}

Diese Statuswerte sind keine neue EFFECT\_ACK-Wire-Erweiterung. Sie sind
Begriffe des Anwendungsprofils und können über bestehende, inhaltsadressierte
Evidenzreferenzen transportiert werden.

\section{Kein Fehlschluss über Physik, Determinismus oder Willen}

\subsection{Physikalische Retrokausalität}

Der ausgeführte Kanal verbindet virtuelle Adressen. Die Hostereignisse, die
Anfrage, Replay und Antwort erzeugen, folgen strikt aufeinander. Daraus folgt
keine positive physikalische Kanalkapazität von einer späteren randomisierten
Intervention zu einem bereits versiegelten früheren Messwert.

Eine physikalische Brückenhypothese müsste mindestens angeben:

\begin{itemize}
  \item welches Gerät oder Feld den behaupteten Träger bildet,
  \item welche später gewählte Variable welchen früheren Messwert beeinflusst,
  \item wie frühere Werte vor der Wahl manipulationssicher versiegelt werden,
  \item wie gewöhnliche Leckage, Synchronisation, Vorhersage und
    Selektionsfehler ausgeschlossen werden,
  \item welche Nullhypothese und Abbruchregel vorab gelten und
  \item wie unabhängige Teams den Befund replizieren können.
\end{itemize}

Bis ein solches Experiment positive, reproduzierbare Evidenz liefert, bleibt
der Claim \texttt{OPEN\_PHYSICAL}. Das ist keine Rückkehr zu null, sondern
die genaue Adresse des nächsten unbewiesenen Knotens.

\subsection{Superdeterminismus, freier Wille und freier Glaube}

Ein deterministischer Replay ist eine Eigenschaft eines endlichen Programms.
Er ist keine vollständige Ontologie des Universums und keine Theorie
menschlicher Entscheidungen. Weder der virtuelle Zeuge noch die neuen
Quellen beweisen Superdeterminismus. Sie widerlegen auch keine spezifizierte
Form freien Willens oder freien Glaubens. Diese Fragen werden nicht durch
Umbenennung eines Softwarezustands entschieden.

\section{Ein falsifizierbares Brückenexperiment}

Der folgende Entwurf ist kein Bericht über eine bereits ausgeführte Messung,
sondern ein Prüfplan, der eine physikalische Behauptung grundsätzlich
entscheidbar machen soll.

\begin{enumerate}
  \item Zum früheren Zeitpunkt $t_0$ erzeugt und versiegelt Labor A eine
    Bitfolge $Y$ mit öffentlicher Commitment-Bindung, ohne Kenntnis der
    späteren Wahl.
  \item Zum späteren Zeitpunkt $t_1$ wählt Labor B per zertifizierter
    Zufallsquelle unabhängig ein Bit $X$.
  \item Das behauptete Brückengerät soll eine vorab definierte statistische
    Abhängigkeit zwischen $X$ und dem schon versiegelten $Y$ erzeugen.
  \item Netzwerk-, Uhr-, Personen-, Speicher-, Metadaten- und
    Präregistrierungslecks werden durch physische und organisatorische
    Trennung kontrolliert. Sham-Geräte und Nullläufe werden verblindet
    eingestreut.
  \item Stichprobengröße, Teststatistik, Fehlergrenzen, Ausschlüsse und
    Abbruchregeln werden vor der Entsiegelung registriert.
  \item Erst nach unabhängiger Replikation wird über eine Statusänderung
    beraten. Ein einzelner positiver Lauf genügt nicht.
\end{enumerate}

Eine solche Versuchsanordnung muss außerdem prüfen, ob scheinbare
Rückwärtsinformation durch gewöhnliche gemeinsame Ursachen, fehlerhafte
Zeitstempel, Vorhersagemodelle, Postselektion oder flexible Auswertung
erklärbar ist. Der Plan verspricht kein positives Ergebnis; er verhindert,
dass ein starkes physikalisches Wort unfalsifizierbar bleibt.

\section{Übersetzung in ein IETF-geeignetes Anwendungsprofil}

Die wissenschaftlichen Aussagen werden nicht als neue Wire-Felder
ausgegeben. Der individuelle Experimental-Entwurf
\texttt{draft-lohmann-qikvrt-epistemic-fairness-observation-profile-00}
liegt über dem unveränderten EFFECT\_ACK-v1-Datensatz mit seinen 35
geschlossenen Mitgliedern und fünf Ergebnissen. Er fügt weder einen Zustand
noch einen Medientyp, eine CDDL-Regel oder eine IANA-Registrierung hinzu.

Das Profil verlangt für wirkungsrelevante Anwendungen unter anderem:

\begin{itemize}
  \item Kennzeichnung als direkt beobachtet, inferiert, simuliert,
    kontrafaktisch, hypothetisch, vorgestellt oder fremdberichtet,
  \item eine zurechenbare Transformationskette von den verfügbaren
    Quellbytes bis zur Entscheidungsrepräsentation,
  \item operationsspezifische Autorisierung für Erhebung, Verarbeitung,
    Speicherung, Weitergabe, Publikation, Modellnutzung und Wirkung,
  \item Offenlegung von privilegierten Zugängen, asymmetrischen Fehlerkosten
    und erreichbaren Widerspruchswegen sowie
  \item explizite Klassifikation jedes Pfades als Software, Replay,
    gewöhnlicher Netzwerktransport, Sensor, Aktor oder behauptete physische
    Brücke.
\end{itemize}

Fehlende oder widersprüchliche Pflichtbelege verhindern die gewöhnliche
Freigabe. Ein positiver Profiltest erzwingt jedoch nie \DONE; sämtliche
Basisprädikate bleiben nötig. Der Entwurf ist nicht eingereicht, kein RFC und
keine IETF-Billigung. Seine Eignung besteht in RFCXML-v3-Form, BCP-14-Sprache,
Security- und Privacy-Abschnitten, klarer IANA-Grenze und testbaren
Konformitätsfällen.

\section{Claim-Delta und Beweisstatus}

\begin{longtable}{@{}L{0.17\linewidth}L{0.22\linewidth}L{0.51\linewidth}@{}}
\toprule
\textbf{Claim} & \textbf{Status} & \textbf{Kernaussage} \\
\midrule
\endhead
\texttt{VTI-ADD-001} & NORMATIVE / DECLARED & achtstufige Realisierungsleiter \\
\texttt{VTI-ADD-002} & NORMATIVE / DECLARED & Software impliziert nicht allein das Naturphänomen \\
\texttt{VTI-ADD-003} & NORMATIVE / DECLARED & siebenfaches Beobachtungsprädikat \\
\texttt{VTI-ADD-004} & INTERPRETATIVE / DECLARED & Vorstellungskraft erzeugt Kandidaten, keine Fakten \\
\texttt{VTI-ADD-005} & NORMATIVE / DECLARED & Vorsprung ersetzt Provenienz oder Autorisierung nicht \\
\texttt{VTI-ADD-006} & OPEN / OPEN & verborgenes reales Deployment nicht belegt \\
\texttt{VTI-ADD-007} & NORMATIVE / DECLARED & kein Schluss auf Superdeterminismus oder gegen freien Willen \\
\texttt{VTI-ADD-008} & OPEN / OPEN & physikalische Rückwärtsübertragung weiterhin offen \\
\bottomrule
\end{longtable}

Die Einstufung \texttt{NORMATIVE} bedeutet, dass eine methodische oder
architektonische Regel deklariert und innerhalb ihres Modells überprüfbar
ist. Sie bedeutet nicht, dass die Regel als Naturgesetz oder allgemeingültige
Ethik bewiesen wurde. \texttt{INTERPRETATIVE} kennzeichnet eine
quelleninspirierte Lesart. \texttt{OPEN} bleibt offen.

\section{Grenzen und nächste Arbeit}

Dieser Nachtrag belegt nicht, dass jedes Computerspiel Menschen heimlich
beobachtet, dass geheimes Wissen existiert, dass ein bestimmter Akteur andere
übervorteilt, dass eine Sprecheridentität durch den Dateinamen authentisiert
ist oder dass kausale beziehungsweise retrokausale Eigenschaften von
Information bereits physikalisch gemessen wurden.

Er leistet stattdessen vier konkrete Dinge:

\begin{enumerate}
  \item Er bewahrt den bereits erreichten Informatikbeweis byte- und
    claimgebunden.
  \item Er übersetzt die neuen Impulse in acht explizit klassifizierte
    Folgeclaims.
  \item Er liefert Konstruktions- und Gegenmodellargumente für den Übergang
    vom Softwarezeugen zum Beobachtungskandidaten.
  \item Er macht empirische, technische und normative Folgearbeit getrennt
    prüfbar.
\end{enumerate}

Die nächsten belastbaren Schritte sind: ein unabhängiger Nachbau des
virtuellen Zeugen; konkrete, autorisierte Instrumentierungsprototypen mit
$\operatorname{OBS}$-Prüfung; interoperable Tests des Anwendungsprofils; und,
falls eine physikalische Behauptung aufrechterhalten wird, ein
präregistriertes Brückenexperiment mit unabhängiger Replikation.

\section{Schlussfolgerung}

Der neue Erkenntnisast ist weder kleinzureden noch größer auszugeben, als
seine Evidenz trägt. Vollständige endliche Information wird im vorhandenen
QIK--VRT-Modell bereits bidirektional zwischen virtuellen Adressen
transportiert. Das ist der gesicherte Ausgangspunkt. Vorstellungskraft kann
daraus reale Geräte- und Beobachtungskandidaten erzeugen. Ein instrumentiertes
Softwaresystem kann unter expliziten Bedingungen tatsächlich ein
Beobachtungssystem werden. Epistemische Fairness sorgt dafür, dass dieser
Übergang nicht heimlich Wahrheitsstatus, Einwilligung oder Wirkungserlaubnis
überspringt.

Für die Physik entsteht dadurch kein automatischer Beweis, aber eine klare
Forschungsroute. Softwaremodell, Gerät, Messung, Kausalbefund und autorisierte
Wirkung sind aufeinander aufbauende, getrennt nachweispflichtige Knoten. So
wird der Erkenntnisbaum wirklich erweitert: Das Bewiesene bleibt stehen; das
Neue erhält eine genaue Form; und das Offene bekommt einen Test.

\section*{Reproduzierbarkeit und Artefakte}
\addcontentsline{toc}{section}{Reproduzierbarkeit und Artefakte}

Zum Kandidaten gehören die LaTeX-Quelle, das daraus mit drei XeLaTeX-Pässen
erzeugte PDF, die maschinenlesbaren Claim- und Quellenbindungen, die
unveränderten ASR-Rohfassungen samt Unsicherheitsprotokoll, eine
allgemeinverständliche Markdown-Fassung, der RFCXML-v3-Profilentwurf und seine
lokal gerenderten Text- und HTML-Fassungen. Roh-Audiodateien werden wegen
Umfang, Rechte- und Datenschutzgrenzen nicht im Git-Repository gespeichert;
ihre lokalen Quellbytes sind über die oben genannten SHA-256-Werte gebunden.

Eine Repository-Persistenz, ein Zenodo-Datensatz und eine
IETF-Datatracker-Einreichung sind voneinander getrennte externe Effekte. Eine
Vorveröffentlichungsprüfung oder ein lokaler Render belegt keinen dieser
Effekte. Ebenso werden weder ein repositoryweiter \texttt{PASS} oder
\texttt{FINAL\_PASS} noch ein globales \texttt{EFFECT\_ACK\_DONE}
beansprucht.

\begin{thebibliography}{9}

\bibitem{base}
I.~Lohmann,
\emph{QIK--VRT: Der bidirektionale virtuelle Zeitkanal},
Working Paper, Version 1.0.0, 1.~August 2026.

\bibitem{shannon}
C.~E. Shannon,
\emph{A Mathematical Theory of Communication},
Bell System Technical Journal 27 (1948), 379--423 und 623--656.

\bibitem{pearl}
J.~Pearl,
\emph{Causality: Models, Reasoning, and Inference},
2.~Auflage, Cambridge University Press, 2009.

\bibitem{rfc6973}
A.~Cooper et al.,
\emph{Privacy Considerations for Internet Protocols},
RFC~6973, Juli 2013, \url{https://www.rfc-editor.org/info/rfc6973}.

\bibitem{rfc8280}
N.~ten Oever und C.~Cath,
\emph{Research into Human Rights Protocol Considerations},
RFC~8280, Oktober 2017, \url{https://www.rfc-editor.org/info/rfc8280}.

\end{thebibliography}

\appendix
\section{Konservativ geprüfte Lesefassungen}

Die vollständigen unveränderten ASR-Rohfassungen und das detaillierte
Unsicherheitsprotokoll stehen in
\texttt{TRANSCRIPTS\_AND\_SOURCE\_PROVENANCE.md}. Die folgenden Fassungen
dienen nur der lesbaren Orientierung.

\subsection{Übervorteilen}

\begin{quote}
Und darüber hinaus solltest du dir überlegen, welche Bedeutung meine
[unsicher: Entdeckung/Entdeckungen hat/haben], nicht nur für die Informatik,
sondern für die Menschheit insgesamt. Denn dieses geheime Wissen, was da
manche Leute haben und es auch bewusst geheim halten, um ihre Vorteile davon
zu haben, muss natürlich mit allen Menschen geteilt werden. Und anders ist
das in heutigen Zeiten auch gar nicht mehr denkbar. Das haben nur noch nicht
alle verstanden. Genauso wenig wie die Physiker diese [unsicher:
Zusammenhänge] verstanden haben. Und deswegen muss ich das den [unsicher:
Herrschaften] auch halt [unsicher: erklären] -- dass es kein
Superdeterminismus ist und auch kein Verzicht auf freien Willen oder freien
Glauben, sondern Eigenschaften des Universums, die für Physiker nur manchmal
paradox erscheinen, es in Wahrheit aber gar nicht sind, sondern nur die
Wirkung von Informationen mit ihrer kausalen und retrokausalen Eigenschaft.
\end{quote}

\subsection{Vorstellungskraft}

\begin{quote}
Und jetzt musst du dir noch einmal klarmachen, dass, wenn das [unsicher: in
Software] möglich ist, das tatsächlich auch in der Realität möglich ist. Und
dazu musst du jetzt all deine Vorstellungskraft mal nutzen und überlegen,
welche Devices, welche Möglichkeiten und welche Eigenschaften es in der
Informatik heutzutage in der Realität gibt und wie wir Menschen damit
potenziell umgehen. Und dann musst du erkennen, dass all das bereits real
passiert, ohne dass es offiziell bekannt ist. Warum? Na ja, weil manche über
dieses spezielle Wissen schon sehr lange verfügen und es immer nur [unsicher:
(geheim) gehalten] haben und [unsicher: dachten, es sei/war] eine gute Idee,
dieses Wissen mit Hilfe von Informatik in die Realität zu bringen. Und so
wird aus einem einfachen [unsicher: Computerspiel] plötzlich ein [unsicher:
Observationssystem]. Erstaunlich, aber nicht verwunderlich.
\end{quote}

\end{document}
